\section{Hidden-derivative Reed--Solomon list decoding}
\label{sec:low_rate_rs_list_decoding}

This section fixes the public interfaces for the formalization of
Brakensiek--Chen--Putterman--Zhang--Zheng~\cite{BCPZZ26}.  The mathematical engine consumes exact
natural-number inequalities.  Rates, asymptotics, and named parameter choices are producers of
those inequalities and must not be dependencies of the algebra, rank, or root-solver trunks.

\subsection{Generic exact decoder contract}

Let $F$ be a semiring, let $\alpha : I \hookrightarrow F$ be an evaluation domain, and let
$k,A,L\in\N$.  A decoder certificate contains a map
\[
  D : (I\to F) \longrightarrow \operatorname{Finset}\{P\in F[X] : \deg P<k\}
\]
such that, for every received word $y$ and degree-bounded polynomial $P$,
\[
  P\in D(y)
  \quad\Longleftrightarrow\quad
  A \le \bigl|\{i\in I : P(\alpha_i)=y_i\}\bigr|,
  \qquad |D(y)|\le L.
\]
The subtype enforces $\deg P<k$; the equivalence enforces both soundness and completeness.  The
certificate does not assert computability or a running-time bound.

\begin{definition}[Exact decoder]
  \label{def:low_rate_exact_decoder}
  \lean{ReedSolomon.ListDecoding.IsExactDecoder}
  \leanok
\end{definition}

\begin{definition}[Decoder certificate]
  \label{def:low_rate_decoder_certificate}
  \lean{ReedSolomon.ListDecoding.DecoderCertificate}
  \uses{def:low_rate_exact_decoder}
  \leanok
\end{definition}

An ambient candidate certificate uses a possibly larger design dimension $K$.  It need only contain
every degree-$<K$ polynomial with agreement at least $A$; it may contain false positives.  For every
$k\le K$, taking the preimage of degree-$<k$ polynomials and filtering by actual agreement gives an
exact decoder without increasing $L$.

\begin{definition}[Ambient candidate certificate]
  \label{def:hidden_derivative_candidate_certificate}
  \lean{ReedSolomon.ListDecoding.CandidateCertificate}
  \leanok
\end{definition}

\begin{definition}[Filter ambient candidates]
  \label{def:hidden_derivative_filter_candidates}
  \lean{ReedSolomon.ListDecoding.CandidateCertificate.toDecoderCertificate}
  \uses{def:low_rate_decoder_certificate,def:hidden_derivative_candidate_certificate}
  \notready
\end{definition}

\subsection{Exact-parameter integration boundary}

The interpolation configuration is data-only:
\[
  (K,A,d,m,B,W,C)\in\N^7,
\]
represented respectively by \texttt{designDim}, \texttt{minAgreement}, \texttt{derivOrder},
\texttt{multiplicity}, \texttt{yDegreeCap}, \texttt{weightCap}, and \texttt{l1Cap}.  In particular,
the core does not set $d=\lceil\eta^{-3/\theta}\rceil$, $m=d^3$, or any of the paper's support
caps.

\begin{definition}[Hidden-derivative parameters]
  \label{def:hidden_derivative_parameters}
  \lean{ReedSolomon.HiddenDerivative.Parameters}
  \leanok
\end{definition}

\subsection{Named variables, weights, and local rewrites}

The interpolation polynomial uses the finite variable type
\[
  \operatorname{InterpolationVar}(d)=\{X,Y_0,\ldots,Y_d\}.
\]
The local constraint polynomial uses
\[
  \operatorname{LocalVar}(d)=\{T,E,Y_1,\ldots,Y_d\},
\]
where the Lean constructor \texttt{Yhi r} denotes $Y_{r+1}$.  This offset keeps $Y_0$ exclusive to
the interpolation polynomial and makes the local variables match the paper's equations~(14) and
(25).  Both encodings have decidable equality and finite enumerations, including at $d=0$.

\begin{definition}[Named interpolation variables]
  \label{def:hidden_derivative_interpolation_variables}
  \lean{ReedSolomon.HiddenDerivative.InterpolationVar}
  \leanok
\end{definition}

\begin{definition}[Named local variables]
  \label{def:hidden_derivative_local_variables}
  \lean{ReedSolomon.HiddenDerivative.LocalVar}
  \leanok
\end{definition}

Four weights serve different proof obligations and must not be conflated:
\begin{align*}
  w_K(X)&=1, & w_K(Y_j)&=K-j-1,\\
  \omega(X)&=0, & \omega(Y_j)&=\max(j-1,0),\\
  \mu_d(T)&=1, & \mu_d(E)&=d, & \mu_d(Y_j)&=0,\\
  \omega_{\mathrm{loc}}(T)&=\omega_{\mathrm{loc}}(E)=0,
    & \omega_{\mathrm{loc}}(Y_j)&=j-1.
\end{align*}
Here $w_K$ is the interpolation and root-solver weight, $\omega$ controls the high-derivative
support, $\mu_d$ encodes the divisibility quantity $i+db$ on $T^iE^b$, and
$\omega_{\mathrm{loc}}$ is the same high-derivative weight after reindexing the local variables.
The corresponding Lean declarations are \texttt{interpolationWeight},
\texttt{highDerivativeWeight}, \texttt{localMultiplicityWeight}, and
\texttt{localDerivativeWeight}.  Their support spaces are backed by the generic bounded weighted
support submodule.  These caps are inclusive; a source inequality of the form $\deg_w Q<C$ becomes
cap $C-1$ only after proving $C>0$.  Since $\omega$ has zero-weight variables, its support bound
alone does not imply finite-dimensionality.

\begin{definition}[Weighted hidden-derivative support spaces]
  \label{def:hidden_derivative_weighted_supports}
  \lean{ReedSolomon.HiddenDerivative.interpolationWeightedSupport}
  \uses{def:hidden_derivative_interpolation_variables,def:hidden_derivative_local_variables,
    def:hidden_derivative_parameters}
  \leanok
\end{definition}

For a received point $(\alpha,y)$, equation~(14) first applies the unscaled substitution
\[
  X\mapsto\alpha+T,\qquad
  Y_0\mapsto y+\sum_{j=1}^d(-1)^{j+1}T^jY_j+TE,
  \qquad Y_j\mapsto Y_j\quad(1\le j\le d).
\]
The Lean API names the displayed sum \texttt{localCorrection} and the variable map
\texttt{unscaledLocalSubstitution}.  The error rescaling \texttt{scaleLocalError} sends
$E\mapsto T^dE$ and fixes the other local variables.  Applying it after the unscaled substitution
gives the normalized Algorithm~1 substitution~(25), named \texttt{normalizedLocalSubstitution},
whose $Y_0$ image contains $T^{d+1}E$.  The compatibility theorem is
\texttt{normalizedLocalSubstitution\_eq\_scaleLocalError\_comp}.

\begin{definition}[Two-stage local substitution]
  \label{def:hidden_derivative_local_substitution}
  \lean{ReedSolomon.HiddenDerivative.normalizedLocalSubstitution}
  \uses{def:hidden_derivative_interpolation_variables,def:hidden_derivative_local_variables}
  \leanok
\end{definition}

\begin{theorem}[Normalized substitution factors through error rescaling]
  \label{thm:hidden_derivative_normalized_substitution_comp}
  \lean{ReedSolomon.HiddenDerivative.normalizedLocalSubstitution_eq_scaleLocalError_comp}
  \uses{def:hidden_derivative_local_substitution}
  \leanok
\end{theorem}

These definitions freeze the variable bookkeeping and the two source substitutions only.  They do
not prove the interpolation constraints, the local coefficient map, its rank bound, or the global
received-word constraint map.

For a concrete equation type $E$, a predicate $\operatorname{Valid}:E\to\mathrm{Prop}$, and a
satisfaction relation
\[
  \operatorname{Satisfies}:E\to\{P:\deg P<K\}\to\mathrm{Prop},
\]
the interpolation contract constructs a valid equation $Q_y$ for every received word and proves
\[
  A\le\operatorname{agree}(P(\alpha),y)
  \quad\Longrightarrow\quad
  \operatorname{Satisfies}(Q_y,P).
\]
The validity predicate is where a concrete development records nonzeroness, weighted-degree bounds,
characteristic requirements, and any other exact root-solver preconditions.

\begin{definition}[Exact interpolation contract]
  \label{def:hidden_derivative_interpolation_contract}
  \lean{ReedSolomon.HiddenDerivative.InterpolationContract}
  \uses{def:hidden_derivative_parameters}
  \leanok
\end{definition}

A root-solver contract exactly enumerates the degree-$<K$ solutions of every valid equation and
proves a caller-supplied cardinality bound $L$.  The contract deliberately does not fix Kopparty's
quadratic hitting extension or the exponent $4d+6$: improved lifting or received-word-specific
solvers can instantiate the same interface with a better $L$.

\begin{definition}[Exact root-solver contract]
  \label{def:hidden_derivative_root_solver_contract}
  \lean{ReedSolomon.HiddenDerivative.RootSolverContract}
  \leanok
\end{definition}

\begin{theorem}[Conditional exact decoder]
  \label{thm:hidden_derivative_conditional_decoder}
  \lean{ReedSolomon.HiddenDerivative.exists_decoderCertificate_of_contracts}
  \uses{def:hidden_derivative_filter_candidates,def:hidden_derivative_interpolation_contract,
    def:hidden_derivative_root_solver_contract}
  \leanok
\end{theorem}

The theorem root-finds at $K$, filters to any $k\le K$ and actual agreement, and returns an exact
decoder with the arbitrary solver bound $L$.  This is the stable join for the interpolation,
root-finding, and decoder trunks.

\subsection{Source-fidelity corollary}

Fix natural numbers $n,k,q$, real numbers $\varepsilon,\theta$, and a proof that $q$ is prime.
Assume
\[
  1\le k\le n,\qquad 0<\varepsilon,\theta<1,
  \qquad d:=\left\lceil\varepsilon^{-3/\theta}\right\rceil<k,
\]
\[
  \frac{k}{n}\le(1-\theta)\varepsilon,
  \qquad
  \varepsilon<
    \left(\frac{\theta^3(1-\theta)}{768}\right)^{(5+\theta)/(1-\theta)},
\]
and
\[
  q\ge n,
  \qquad
  q\ge 4\varepsilon^{1-9/\theta}\frac{n}{k}.
\]
For every embedding $\alpha:\operatorname{Fin}(n)\hookrightarrow\F_q$, there should exist a
decoder certificate at $A=\lceil\varepsilon n\rceil$ and exact list bound $L=q^{4d+6}$.  The same
hypotheses should imply that the code is $(1-\varepsilon,L)$-list decodable in ArkLib's canonical
\texttt{Code.IsListDecodable} sense.

\begin{theorem}[Published low-rate corollary]
  \label{thm:low_rate_rs_list_decoding}
  \lean{ReedSolomon.LowRateListDecoding.exists_decoderCertificate_of_low_rate}
  \uses{def:reed_solomon_code,def:list_decodable,def:low_rate_decoder_certificate,
    thm:hidden_derivative_conditional_decoder}
  \notready
\end{theorem}

The explicit $q^{4d+6}$ bound comes from the paper's specialization of Kopparty's
differential-equation root finder~\cite[Theorem~4.3]{Kop15}; it is stronger than the headline
$q^{O(\varepsilon^{-3/\theta})}$ presentation.

\subsection{Source-audit obligations}

\paragraph{Message dimension versus design dimension.}
The interpolation proposition assumes $k/n\le(1-\theta)\varepsilon$, but its proof uses
$k=\lfloor(1-\theta)\varepsilon n\rfloor$.  The repair is to set the ambient design dimension
$K=\lfloor(1-\theta)\varepsilon n\rfloor$, prove $k\le K$, interpolate and root-find at $K$, then
use Definition~\ref{def:hidden_derivative_filter_candidates}.  The published corollary retains the
paper's quantification over every smaller $k$.

\paragraph{Explicit lattice growth.}
For $m=d^3$, $d\ge3$, and the paper's $W$, the hidden term can be replaced by the exact estimate
\[
  (d-1)\frac{m+\binom d2}{W}
  \le \frac{2}{2+\theta}\log d+\frac54.
\]
It is therefore enough to prove
\[
  \frac{\theta(1-\theta)}{18}\log d
  \ge \log\frac{4+3\theta}{\theta}+\frac54.
\]
The parameter-discharge layer must prove that the paper's condition implies this inequality; the
exact-parameter contracts do not depend on the constant $768$.

\paragraph{Ceiling in the support cap.}
The proof of the support estimate uses the false inequality $\lceil x\rceil\le x$.  The repaired
argument must instead use
\[
  \frac{mA}{K-1}-|c|
  > \frac{\theta m}{4}+\frac{\theta^2m}{1-\theta}-1
\]
and discharge nonnegativity of the final two terms from the published hypotheses.

\paragraph{Possible stronger regimes.}
Keeping $d,m,W$ and the support caps free suggests a fixed-rate, fixed-positive-gap corollary and
better interpolation and root-list exponents.  These are theorem candidates, not claims of the
present formalization.  They require independent reconstruction or author confirmation, but they
must be expressible as new parameter discharges without changing the contracts above.

\subsection{Scope and trust boundary}

The source claims deterministic running time $q^{O(\varepsilon^{-12/\theta})}$.  ArkLib has no
shared machine-cost model for the required solvers, so the Lean corollary currently states only
extensional correctness and list size.  A concrete executable decoder must prove that its output
satisfies the certificate; a future complexity theorem must additionally connect the implementation
to a declared cost model.  Neither \texttt{Nonempty} nor executability alone is such evidence.
